% =====================================================================
%  SKILL 2026 SHORT PAPER  (4-6 pages incl. references)
%
%  PORTING NOTE (do this before submission):
%   - Replace the preamble below with the official GI-LNI template:
%         \documentclass[english]{lni}
%     and use the LNI front-matter macros (\title, \author{...}{...}{...},
%     \maketitle, \keywords{...}) from the template instead of the
%     article-class title block used here for preview.
%   - This file compiles as-is with pdflatex (article class) so you can
%     read it; the body (sections, math, figure, thebibliography) ports
%     verbatim into the LNI class.
%   - VERIFY EVERY BIBLIOGRAPHY ENTRY. arXiv IDs marked %VERIFY were taken
%     from search results in late development and must be checked against
%     the canonical record, especially the 2026 preprints.
%   - Fill in author name, affiliation, e-mail, and acknowledgements.
%   - The figure file grokking_normal_hyperbolicity.png must sit next to
%     this .tex file.
% =====================================================================
\documentclass[11pt]{article}
\usepackage[utf8]{inputenc}
\usepackage[T1]{fontenc}
\usepackage[margin=1in]{geometry}
\usepackage{amsmath,amssymb}
\usepackage{graphicx}
\usepackage{booktabs}
\usepackage{hyperref}
\usepackage{microtype}

\title{Is Grokking a Loss of Normal Hyperbolicity\\ of the Interpolation Manifold?}
\author{Suvinava [Last Name]\thanks{[Affiliation], [e-mail]. Correspondence welcome.}}
\date{}

\newcommand{\sminp}{\sigma_{\min}^{+}}
\newcommand{\R}{\mathbb{R}}

\begin{document}
\maketitle

\begin{abstract}
A recent line of work recasts the post-memorization phase of grokking as
constrained optimization: once a network interpolates the training set,
weight decay drives a slow drift along the zero-loss manifold toward lower
norm. This is, in the language of dynamical systems, a fast--slow system in
which the interpolation manifold plays the role of a slow manifold. We ask a
question that this framing makes natural but that the existing literature
does not address: \emph{is the sharp generalization transition a loss of
normal hyperbolicity of that manifold} -- a fold- or bifurcation-like event
in which a normal restoring direction goes flat -- or does the manifold stay
uniformly attracting while generalization happens by smooth drift? We propose
a simple, optimizer-agnostic diagnostic: the smallest nonzero singular value
$\sminp(J)$ of the residual Jacobian, which equals the smallest normal
restoring curvature of the manifold. On a two-layer ReLU network trained to
grok modular addition under squared loss, $\sminp(J)$ does not collapse at the
transition; it is near zero only \emph{before} memorization and attains its
largest values \emph{during} the transition. This is preliminary evidence
against the bifurcation hypothesis and in favor of the smooth-contraction
picture. We are explicit that a single-seed, single-setting, gradual-transition
experiment under Adam does not prove the absence of a bifurcation; it
constrains where one could hide.
\end{abstract}

\section{Introduction}
Grokking is the empirical observation that a network can reach perfect
training accuracy with near-chance test accuracy and only generalize after a
long additional period of training~\cite{power2022}. A productive recent
account treats the long post-memorization phase as constrained optimization.
\citet{musat2025} argue that, after interpolation, gradient descent with small
weight decay effectively minimizes the parameter norm subject to staying on
the zero-loss manifold, proving this in the limit of infinitesimal learning
rate and weight decay, and a concurrent framework reaches the same endpoint
through a Riemannian norm flow on the manifold of interpolating
solutions~\cite{boursier2025}. \citet{lyu2024} prove, for homogeneous networks
with large initialization and small weight decay, a sharp transition from a
kernel predictor to a KKT point of a global min-norm/max-margin problem on
both classification and regression tasks, and identify modular addition under
cross-entropy as the motivating phenomenon.

All of these accounts share a structural picture: a fast process pulls the
parameters onto a low-dimensional set of (near-)interpolating solutions, and a
slow process driven by weight decay moves the parameters \emph{along} that
set. This is exactly the structure of a singularly perturbed, or fast--slow,
dynamical system, with the interpolation manifold as the slow manifold. Yet
none of these works invokes the corresponding mathematics -- normal
hyperbolicity, Fenichel persistence of slow manifolds~\cite{fenichel1979,
kuehn2015}, and what happens when normal hyperbolicity is \emph{lost}. The
question we raise is whether grokking's defining feature, the suddenness of
the transition, is a signature of that loss.

Concretely, we contrast two hypotheses.
\textbf{(H1)} The interpolation manifold loses normal hyperbolicity at the
transition: a normal restoring direction softens toward zero, the
timescale separation between fast and slow motion breaks down, and the
trajectory makes a fast jump -- a fold or pitchfork in the fast--slow
decomposition. Modular addition's symmetry group makes a symmetry-breaking
pitchfork a natural candidate normal form.
\textbf{(H0)} The manifold stays uniformly normally hyperbolic; the transition
is a feature of the slow drift through a region where test behavior changes
quickly, without any degeneration of the manifold's geometry. This is the
picture implied by the contraction-based accounts above.

H1 is the more interesting hypothesis and, if true, would tie the dynamical
and mechanistic-interpretability views of grokking together. The contribution
of this short paper is a clean test that distinguishes H1 from H0 and a
preliminary -- and, as it turns out, negative -- result.

\section{The diagnostic}
Let $f(\theta;x)\in\R^{c}$ denote the network output and consider the squared
loss $L(\theta)=\frac1n\sum_{i=1}^{n}\lVert f(\theta;x_i)-y_i\rVert^2$ on a
training set of size $n$. We use squared loss precisely because it makes the
\emph{interpolation manifold} $\mathcal{M}=\{\theta : f(\theta;x_i)=y_i\ \forall i\}$
an exact object with a well-defined tangent and normal structure, which the
diagnostic needs. Training with decoupled weight decay $\lambda$ follows, in
the gradient-flow limit, $\dot\theta=-\nabla L(\theta)-\lambda\theta$; for
small $\lambda$ the term $-\nabla L$ is fast and pulls toward $\mathcal{M}$,
while $-\lambda\theta$ is the slow drift along it.

Collect the residuals into $r(\theta)=\big(f(\theta;x_i)-y_i\big)_i\in\R^{nc}$
and let $J(\theta)=\partial r/\partial\theta\in\R^{nc\times D}$ be the residual
Jacobian. The Gauss--Newton Hessian is $G(\theta)=J(\theta)^\top J(\theta)$. At
a point of $\mathcal{M}$, the row space of $J$ is the normal space
$N_\theta\mathcal{M}$ (directions in which residuals -- hence the loss --
change to first order), and $\ker J$ is the tangent space $T_\theta\mathcal{M}$
(directions along which the model stays interpolating). The nonzero
eigenvalues of $G$ are the squared singular values $\sigma_i(J)^2$, so the
\emph{smallest normal restoring curvature} of the manifold is
\begin{equation}
  \kappa(\theta) \;=\; \sminp\!\big(J(\theta)\big)^2 ,
\end{equation}
where $\sminp$ denotes the smallest nonzero singular value. In an
over-parameterized network with $nc<D$, $J$ is generically of full row rank, so
all $nc$ singular values are positive and $\sminp$ is simply the smallest of
them.

Normal hyperbolicity of $\mathcal{M}$ for the fast subsystem is the statement
that $\kappa(\theta)$ is bounded away from zero, uniformly along the slow
drift; this is the condition under which Fenichel theory guarantees that the
true (perturbed) slow manifold persists and the reduced flow is the norm flow
studied by~\cite{musat2025,boursier2025}. A loss of normal hyperbolicity --
the prerequisite for the fold/pitchfork mechanism of H1 -- is exactly
$\sminp(J)\to 0$. The diagnostic is therefore the single scalar trajectory
$\sminp(J(\theta(t)))$, read against the test-accuracy curve. A dip toward zero
at the transition supports H1; a value bounded away from zero supports H0. The
quantity is a property of the model and data at $\theta$, so it is well-defined
for any optimizer, not only the gradient flow for which the slow-drift theory
is cleanest.

\section{Experimental setup}
We train a two-layer ReLU network $f(\theta;x)=W_2\,\mathrm{ReLU}(W_1 x + b_1)$
with width $96$ and no second-layer bias on modular addition mod $p=11$: the
input is the concatenation of the one-hot encodings of $a$ and $b$, and the
target is the one-hot encoding of $(a+b)\bmod p$. We use a random
$70\%/30\%$ train/test split (seed fixed), squared loss to the one-hot targets,
full-batch optimization, and a large initialization (Kaiming scaled by $3.5$)
to induce a clear memorization plateau. To reach the transition within a
practical compute budget we optimize with AdamW (learning rate $3\times10^{-3}$,
decoupled weight decay $2.0$) for $35{,}000$ steps. We snapshot $\theta$ every
$300$ steps and compute $\sminp(J)$ on the snapshots in a separate pass, so the
cost of the singular value decomposition does not slow training. We return to
the choice of Adam, rather than the gradient descent of the underlying theory,
in the discussion.

\section{Results}
The run groks (Figure~\ref{fig:main}): training accuracy reaches $1.0$ by about
step $4{,}000$, test accuracy stays at $0.0$ through about step $6{,}000$, and
then rises over roughly $7{,}000$--$17{,}000$ steps to $\approx 0.95$.

The diagnostic does not behave as H1 predicts. $\sminp(J)$ is near zero only
\emph{before} memorization (its global minimum, $\approx 7\times10^{-3}$, occurs
at step $1{,}200$, while training accuracy is still well below $1$ and the
network has not yet differentiated its outputs). It then rises monotonically as
the interpolation manifold forms, and across the entire generalization
transition it sits at its largest values; if anything, the normal restoring
curvature \emph{peaks} mid-transition rather than collapsing.
Table~\ref{tab:phases} summarizes the medians by phase.

\begin{figure}[t]
  \centering
  \includegraphics[width=0.92\linewidth]{grokking_normal_hyperbolicity.png}
  \caption{Test and training accuracy (left axis) and the normal restoring
  curvature $\sminp(J)$ (right axis) over training. The shaded band marks the
  generalization transition. $\sminp(J)$ does not dip during the transition;
  it is small only before memorization and is largest while the model
  generalizes. Two-layer ReLU network, squared loss, AdamW, modular addition
  $p=11$; single seed.}
  \label{fig:main}
\end{figure}

\begin{table}[t]
  \centering
  \begin{tabular}{lccc}
    \toprule
    Phase & Step range & median $\sminp(J)$ & min $\sminp(J)$\\
    \midrule
    Pre-memorization & $600$--$1800$ & $\approx 0.013$ & $0.0072$\\
    Memorization plateau & $3000$--$6000$ & $0.13$ & $0.045$\\
    Generalization transition & $7000$--$17000$ & $0.22$ & $0.205$\\
    Post-transition & $\ge 22000$ & $0.18$ & $\approx 0.18$\\
    \bottomrule
  \end{tabular}
  \caption{Normal restoring curvature by training phase. The only regime in
  which $\sminp(J)$ approaches zero is pre-memorization, which is unrelated to
  grokking; during the transition it is at its maximum. The pre-memorization
  reading also shows the diagnostic is sensitive enough to register a genuinely
  near-degenerate Jacobian when one occurs.}
  \label{tab:phases}
\end{table}

A secondary observation: the parameter norm does not decrease through the
transition in this run (it stays near $10$ and drifts slightly upward). The
generalization event is therefore not accompanied by visible $\ell_2$-norm
contraction here. This is consistent with the implicit bias of AdamW being
governed by an $\ell_\infty$-type rather than $\ell_2$ geometry~\cite{xie2024,
gronich2026}, and we treat it as a caveat about the optimizer rather than a
claim about the manifold.

\section{Discussion and limitations}
Taken at face value, the experiment favors H0 over H1: the interpolation
manifold remains robustly normally hyperbolic across grokking, and the
suddenness of the transition is not mirrored by any softening of a normal
restoring direction. This is the outcome one expects from the smooth
contraction picture of~\cite{musat2025,boursier2025,lyu2024} and from recent
work framing the grokking delay as exponential norm contraction across a
geometric gap.\footnote{A 2026 preprint derives a norm-separation delay law for
discrete regularized SGD; see the bibliography entry marked \%VERIFY. We cite
it for context and have not independently reproduced it.} Because the
entry--exit delay law of a delayed-loss-of-stability mechanism only applies if
a bifurcation exists, the negative result also removes the motivation for that
particular theoretical route.

We are deliberate about what this does \emph{not} establish. The evidence is
preliminary and bounded in at least five ways. First, it is a single seed and
a single architecture and hyperparameter setting at small $p$; the curve should
be repeated across seeds and moduli. Second, we optimize with AdamW, whereas
the fast--slow theory is stated for gradient descent or gradient flow; the
clean theory regime is untested here. Third, the observed transition is
moderately gradual rather than maximally sharp, and the sharpest grokking
regimes -- under gradient descent with large initialization and small weight
decay -- are the most stringent place to look for a fold. Fourth, $\sminp(J)$
is the \emph{global} smallest normal singular value; a loss of normal
hyperbolicity confined to a low-dimensional, test-relevant subspace could in
principle occur while the global minimum stays bounded away from zero, and a
subspace-resolved version of the diagnostic would be needed to exclude it.
Fifth, we measure a discrete trajectory, not the gradient-flow limit in which
the manifold structure is exactly defined. For these reasons we frame the
result as constraining, not refuting: in this regime, grokking is not
accompanied by a global collapse of the manifold's normal restoring curvature.

These limitations also define the experiments that would settle the question:
the same diagnostic under gradient descent in a sharp-transition regime; a
subspace-resolved curvature that projects onto the directions separating train
from test points; and, if a dip never appears, a positive characterization of
\emph{which} geometric quantity does move monotonically through the transition,
which under AdamW may be an $\ell_\infty$- or sparsity-based measure rather than
$\ell_2$ norm.

\section{Conclusion}
We asked whether the sharp transition in grokking is a loss of normal
hyperbolicity of the interpolation manifold, and we gave a diagnostic --- the
smallest nonzero singular value of the residual Jacobian --- that turns the
question into a single measurable curve. On a network that groks modular
addition under squared loss, the answer appears to be no: the manifold stays
normally hyperbolic through the transition. We do not claim to have proved the
absence of a bifurcation; we have shown that, in a clean grokking run with a
sensitive instrument, no such signature is present, which is itself a
constraint on theories of why the transition is sharp.

\paragraph{Reproducibility.} Training and diagnostic code, and the data behind
Figure~\ref{fig:main}, are available at [repository link].

\begin{thebibliography}{99}\small
\bibitem{power2022} A.~Power, Y.~Burda, H.~Edwards, I.~Babuschkin, V.~Misra.
  Grokking: Generalization beyond overfitting on small algorithmic datasets.
  arXiv:2201.02177, 2022.
\bibitem{musat2025} T.~Mu\c{s}at. The Geometry of Grokking: Norm Minimization
  on the Zero-Loss Manifold. arXiv:2511.01938, 2025.
\bibitem{boursier2025} E.~Boursier et al. A Theoretical Framework for Grokking:
  Interpolation followed by Riemannian Norm Minimisation. arXiv:2505.20172,
  2025. %VERIFY id and author list
\bibitem{lyu2024} K.~Lyu, J.~Jin, Z.~Li, S.~S.~Du, J.~D.~Lee, W.~Hu. Dichotomy
  of Early and Late Phase Implicit Biases Can Provably Induce Grokking. ICLR
  2024; arXiv:2311.18817.
\bibitem{nanda2023} N.~Nanda, L.~Chan, T.~Lieberum, J.~Smith, J.~Steinhardt.
  Progress measures for grokking via mechanistic interpretability. ICLR 2023;
  arXiv:2301.05217.
\bibitem{soudry2018} D.~Soudry, E.~Hoffer, M.~S.~Nacson, S.~Gunasekar,
  N.~Srebro. The implicit bias of gradient descent on separable data. JMLR
  19(1):2822--2878, 2018.
\bibitem{lyuli2020} K.~Lyu, J.~Li. Gradient descent maximizes the margin of
  homogeneous neural networks. ICLR 2020.
\bibitem{xie2024} S.~Xie et al. Implicit bias of AdamW: $\ell_\infty$-norm
  constrained optimization. ICML 2024; arXiv:2404.04454. %VERIFY
\bibitem{gronich2026} A.~Gronich, G.~Vardi. The Implicit Bias of Adam and Muon
  on Smooth Homogeneous Neural Networks. arXiv:2602.16340, 2026. %VERIFY
\bibitem{fenichel1979} N.~Fenichel. Geometric singular perturbation theory for
  ordinary differential equations. J.\ Differential Equations 31(1):53--98,
  1979.
\bibitem{kuehn2015} C.~Kuehn. Multiple Time Scale Dynamics. Springer, 2015.
\bibitem{omnigrok} Z.~Liu, E.~J.~Michaud, M.~Tegmark. Omnigrok: Grokking beyond
  algorithmic data. ICLR 2023. %VERIFY id
\bibitem{delaylaw} The Norm-Separation Delay Law of Grokking.
  arXiv:2603.13331, 2026. %VERIFY id and authors; post-dates our main analysis
\end{thebibliography}

\end{document}
